\documentclass[12pt,a4paper]{article}
\usepackage{ra_preamble}
\title{\textbf{Wave Function Collapse as Irreversible Entropy Production:\\
The Universal Actualization Timescale}\\[0.4em]
\large\normalfont Relational Actualism Paper~II\\
\normalsize Under review, \textit{Foundations of Physics}}
\author{Joshua F.~Sandeman\\[0.3em]
\small Independent Researcher, Salem, Oregon\\
\small \texttt{sansuikyo@gmail.com} $\cdot$ \texttt{relationalactualism.org}}
\date{April 2026}

\begin{document}
\maketitle
\begin{abstract}
We present Relational Actualism's resolution of the quantum measurement problem and derive a universal actualization timescale formula that applies from the Planck scale to biological systems.  Wavefunction collapse is not a postulate but a physical event: an interaction actualizes when $\Delta S(\rho\|\sigma_0) > \DeltaS = 0.60069\,\mathrm{nats}$.  This threshold is derived from the BDG causal set action in $d=4$ (Paper~I \cite{RA_P1}).  We prove causal invariance of the quantum measure for any permutation of spacelike events \statusLV\ (O01, O02) and resolve the Unruh paradox: the Rindler thermal state is KMS, giving $\Delta S_{\mathrm{per\;photon}} = 0$ and $\tact \to \infty$ in Minkowski vacuum \statusDR\ (O05r, from L05 \statusLV).  The universal formula $\tact = \DeltaS / \Gamma_{\mathrm{on\text{-}shell}}$ spans fourteen orders of magnitude: at $T=1\,\mathrm{mK}$ (BMV experiment), $\tact \approx 1.9\times10^4\,\mathrm{s} \gg \tau_{\mathrm{exp}}$; at $T=300\,\mathrm{K}$ (biology), $\tau_d \approx 20\,\mathrm{fs} = \tau_{\mathrm{vib}}$.  The actualization timescale scales as $T^{-4}$ through the blackbody rate $\Gamma_{bb}\propto T^4$, in contrast to the standard Lindblad result $\tau_{\mathrm{dec}} \sim T^{-1}$.  This difference in power law is a falsifiable prediction distinguishing RA from all decoherence-based approaches.  A parameter-free spin-bath collapse prediction $t^* = 0.274/g$ is also derived.
\end{abstract}

\tableofcontents

%───────────────────────────────────────────────────────────────
\section{The Measurement Problem and Its Resolution}

The measurement problem is structural: quantum mechanics describes superpositions, but every observation produces a definite outcome.  Every proposed resolution either adds structure (many worlds, hidden variables), modifies mathematics (GRW, CSL, pilot waves), or refuses engagement (Copenhagen).  None derives collapse from more fundamental principles.

RA dissolves the problem by replacing the Heisenberg cut with an objective criterion: an interaction \emph{actualizes} — becomes a permanent event — when it produces a quantum relative entropy change exceeding $\DeltaS = 0.60069\,\mathrm{nats}$.  The quantum-to-classical transition is not a boundary between systems but a threshold that any interaction either crosses or does not.

\paragraph{The Schr\"odinger cat.} A real cat is a densely self-coupled actualization network: $\sim10^{25}$ atoms exchanging on-shell bosons at rates $\sim10^{13}\,\mathrm{s}^{-1}$ at $300\,\mathrm{K}$.  The decoherence timescale $\tau_d \approx 20\,\mathrm{fs}$ (derived in Section~\ref{sec:tau}) means the cat undergoes $\sim5\times10^{13}$ actualization events per second.  It is classical by constitution, not observation.  The radioactive isotope is the genuine quantum system — a single nuclear degree of freedom with a well-defined (but unpredictable) actualization time.

%───────────────────────────────────────────────────────────────
\section{The Actualization Criterion}

\subsection{Objective Threshold}

An interaction constitutes an actualization event if and only if:
\begin{equation}
  \Delta S(\rho \| \sigma_0) > \DeltaS = 0.60069\,\mathrm{nats}
  \label{eq:criterion}
\end{equation}
where $\sigma_0$ is the vacuum state.  This is observer-independent and Lorentz-covariant:
\begin{leanbox}
\textbf{Theorem (L04, Frame Independence) \statusLV:} $S(U\rho U^\dagger \| \sigma_0) = S(\rho\|\sigma_0)$ for all Poincar\'e $U$ with $U\sigma_0 U^\dagger = \sigma_0$ (\texttt{frame\_independence}, RA\_AQFT\_Proofs\_v10.lean).
\end{leanbox}
In the perturbative regime, Eq.~(\ref{eq:criterion}) reduces to the on-shell condition $p^\mu p_\mu = m^2c^2$.  Off-shell virtual processes ($\Delta S = 0$) never cross the threshold and do not contribute to the actualized causal graph.

\subsection{Causal Invariance (Lean-Verified)}

\begin{leanbox}
\textbf{Theorem (O01, Amplitude Locality) \statusLV:} The BDG amplitude $a(v|C)$ depends only on $C \cap \mathrm{past}(v)$ (\texttt{bdg\_amplitude\_locality}, RA\_AmpLocality.lean, zero sorry tags).

\textbf{Theorem (O02, Causal Invariance) \statusLV:} The quantum measure is independent of the linear extension $\sigma$ for \emph{any} permutation $\sigma.\mathrm{Perm}\,\sigma'$ (\texttt{bdg\_causal\_invariance}, RA\_AmpLocality.lean).
\end{leanbox}

%───────────────────────────────────────────────────────────────
\section{The Universal Actualization Timescale}\label{sec:tau}

\subsection{Derivation}

The actualization timescale is the mean time for a system to accumulate $\Delta S > \DeltaS$ through on-shell boson exchanges with its environment:
\begin{equation}
  \boxed{\tact(m,T) = \frac{\DeltaS}{\Gamma_{\mathrm{on\text{-}shell}}(m,T) \cdot \Delta S_{\mathrm{per\;boson}}(T)}}
  \label{eq:tau}
\end{equation}
Here $\Gamma_{\mathrm{on\text{-}shell}}$ is the rate of on-shell boson exchanges (those satisfying $\hbar\omega > k_BT\cdot\DeltaS$), and $\Delta S_{\mathrm{per\;boson}} \approx \hbar\omega/(k_BT)$ is the entropy carried per exchange at the thermal peak.

\subsection{Comparison with the Lindblad Master Equation}\label{sec:lindblad}

Standard decoherence theory uses the Lindblad master equation:
\begin{equation}
  \dot\rho = -\frac{i}{\hbar}[H,\rho] + \sum_k \Bigl(L_k\rho L_k^\dagger - \tfrac{1}{2}\{L_k^\dagger L_k,\rho\}\Bigr)
  \label{eq:lindblad}
\end{equation}
The decoherence rate $\Gamma_{\mathrm{dec}}$ is computed from the environmental spectral density $J(\omega)$, giving (for the quantum Brownian motion model \cite{Joos1985,Zurek2003}):
\begin{equation}
  \tau_{\mathrm{dec}}^{(\mathrm{Lindblad})} \sim \frac{\hbar}{k_BT}\left(\frac{\lambda_{\mathrm{th}}}{\Delta x}\right)^2 \propto T^{-1}
  \label{eq:lindblad_tau}
\end{equation}
The RA prediction differs in two ways:

\begin{table}[h]
\centering
\caption{RA vs.\ Lindblad: actualization timescale comparison}
\label{tab:comparison}
\begin{tabular}{lll}
\toprule
Property & Lindblad (standard QM) & RA (this work) \\
\midrule
Temperature scaling & $\tact \propto T^{-1}$ & $\tact \propto T^{-4}$ (via $\Gamma_{bb}\propto T^4$) \\
Threshold behavior & Smooth exponential & Step function at $T^*$ \\
Virtual processes & Contribute to decoherence & $\Delta S=0$, do not actualize \\
BMV at $1\,\mathrm{mK}$ & $\tau \sim 10^2\,\mathrm{s}$ & $\tact \sim 10^4\,\mathrm{s}$ \\
Biology at $300\,\mathrm{K}$ & $\tau \sim 10^{-13}\,\mathrm{s}$ & $\tact \sim 2\times10^{-14}\,\mathrm{s}$ \\
\bottomrule
\end{tabular}
\end{table}

\textbf{Where standard theory breaks down:} The Lindblad equation treats the environmental bath as a continuous spectral density extending to $\omega=0$.  RA imposes a hard threshold: modes with $\hbar\omega < k_BT\cdot\DeltaS$ are off-shell and contribute $\Delta S=0$.  At low temperatures ($T\ll \hbar\omega_{\mathrm{vib}}/k_B$), the on-shell mode count drops sharply, replacing the smooth $T^{-1}$ Lindblad scaling with the $T^4$ blackbody suppression.  This distinction is absent in all continuous-bath decoherence approaches.

\textbf{The step-function signature:} RA predicts a critical temperature $T^*$ below which $\tact \gg \tau_{\mathrm{exp}}$ (quantum behavior preserved) and above which $\tact \ll \tau_{\mathrm{exp}}$ (classical).  Standard QM predicts smooth exponential suppression.  The step function is the qualitative experimental signature.

\subsection{The BMV Experiment ($T\approx 1\,\mathrm{mK}$)}

For a levitated nanosphere with $m \approx 10^{-14}\,\mathrm{kg}$, $R \approx 780\,\mathrm{nm}$:
\begin{align}
  \frac{\hbar\omega_{\mathrm{vib}}}{k_BT}\biggr|_{1\,\mathrm{mK}} &\approx 76{,}000 \gg 1 \quad\text{(zero phonons on-shell)} \\
  \Gamma_{bb}(T=1\,\mathrm{mK}) &= \frac{\sigma_{SB} A T^4}{E_{\mathrm{photon}}} \approx 10^{-5}\,\mathrm{s}^{-1} \\
  \tact &= \frac{0.601}{10^{-5}\times 2.82} \approx 1.9\times 10^4\,\mathrm{s} \gg \tau_{\mathrm{exp}} \approx 1\text{--}10\,\mathrm{s}
  \label{eq:bmv}
\end{align}
The center-of-mass remains unactualized throughout the BMV experiment.  Gravity-mediated entanglement is zero \statusDR.  The mechanism: the $T^4$ suppression of blackbody flux at millikelvin temperatures makes the on-shell exchange rate $\Gamma_{bb}$ too small to accumulate $\DeltaS$ on any experimentally relevant timescale.

\subsection{The Biological Scale ($T\approx 300\,\mathrm{K}$)}

At body temperature with optical phonons ($\omega_{\mathrm{opt}} = 3\times10^{13}\,\mathrm{rad/s}$):
\begin{align}
  \frac{\hbar\omega_{\mathrm{opt}}}{k_BT}\biggr|_{300\,\mathrm{K}} &= 0.76 \quad\text{(on-shell: }{\hbar\omega}/{k_BT} > \DeltaS = 0.6) \\
  \Gamma_{\mathrm{opt}} &= \frac{k_BT}{\hbar} = 3.93\times10^{13}\,\mathrm{s}^{-1} \\
  \tau_d &= \frac{0.601}{3.93\times10^{13}\times 0.76} = 2.0\times10^{-14}\,\mathrm{s} = 20\,\mathrm{fs}
  \label{eq:bio}
\end{align}
This equals $\tau_{\mathrm{vib}}$, the vibrational decoherence timescale.  The RA on-shell selection criterion identifies the correct bath (vibrational, not Debye acoustic) without additional input.

%───────────────────────────────────────────────────────────────
\section{The Unruh Paradox Resolved}\label{sec:unruh}

\begin{leanbox}
\textbf{Theorem (L05, Rindler Stationarity) \statusLV:} $\frac{d}{dt}S(\alpha_t(\rho_\beta)\|\sigma_0) = 0$ under modular flow $\alpha_t$ (\texttt{rindler\_stationarity}, RA\_AQFT\_Proofs\_v10.lean).

\textbf{Theorem (L06, Rindler Thermal) \statusLV:} $\rho_\beta$ is a valid density matrix (Hermitian, positive, $\mathrm{tr}=1$) (\texttt{rindlerThermal}).
\end{leanbox}

Resolution via the universal timescale (O05r \statusDR): A Rindler observer perceives the Minkowski vacuum as thermal at $T_U = \hbar a/(2\pi c k_B)$.  L05 establishes that the Rindler bath is KMS — by definition, each photon exchange with a KMS bath satisfies $\Delta S_{\mathrm{rel}} = 0$ per exchange.  Therefore:
\begin{equation}
  \Delta S_{\mathrm{per\;photon}}(\text{Rindler in Mink. vac.}) = 0 \implies \tact = \frac{\DeltaS}{0} \to \infty
\end{equation}
The detector never fires, regardless of acceleration.  The Unruh radiation is real (L05--L06) but produces zero relative entropy change relative to the vacuum because it \emph{is} the vacuum in another frame.

%───────────────────────────────────────────────────────────────
\section{Spin-Bath Collapse Prediction}

For a spin-half system coupled to a bosonic bath with coupling constant $g$ \statusCV:
\begin{equation}
  t^* = \frac{1}{g}\sqrt{\frac{\DeltaS}{2}} = \frac{0.274}{g}
  \label{eq:spinbath}
\end{equation}
This is the time at which accumulated relative entropy first exceeds $\DeltaS$.  The $g$-dependence and numerical prefactor $0.274 = \sqrt{\DeltaS/2}$ are fully determined predictions, distinguishable from GRW ($t^* \propto m^{-1}$, linear in inverse mass) and CSL models (different parameter dependence).

%───────────────────────────────────────────────────────────────
\section{Discussion}

The universal formula~(\ref{eq:tau}) provides a unified account of decoherence across scales with no free parameters.  The key predictions are:
\begin{enumerate}
  \item \textbf{BMV null result} (near-term): $\tact \approx 1.9\times10^4\,\mathrm{s} \gg 1\text{--}10\,\mathrm{s}$ at $T=1\,\mathrm{mK}$.
  \item \textbf{$T^{-4}$ scaling} (distinctive): measurable by varying temperature across the step-function threshold $T^*$.  Standard QM gives $T^{-1}$.
  \item \textbf{Spin-bath}: $t^* = 0.274/g$, parameter-free, testable with superconducting circuits.
  \item \textbf{Biology}: $\tau_d = 20\,\mathrm{fs}$ at $T=300\,\mathrm{K}$, equal to $\tau_{\mathrm{vib}}$.
\end{enumerate}
All four are derivable from a single equation with a single derived constant ($\DeltaS = 0.60069$).

\section*{Acknowledgments}
The universal timescale formula and its connection to Paper~III--IV was developed with Claude (Anthropic).  Lean~4 files at \url{https://github.com/jsandeman/Relational-Actualism}.

\bibliographystyle{plain}
\begin{thebibliography}{99}
\bibitem{RA_P1} Sandeman, J.~F. (2026). Paper~I: Coupling Constants. Zenodo, doi:10.5281/zenodo.19362289.
\bibitem{RA_P5} Sandeman, J.~F. (2026). Paper~V: The Complete Framework. Zenodo, doi:10.5281/zenodo.19363077.
\bibitem{Bell1987} Bell, J.~S. (1987). \textit{Speakable and Unspeakable in Quantum Mechanics}. Cambridge.
\bibitem{GRW1986} Ghirardi, G.~C., Rimini, A., and Weber, T. (1986). Unified dynamics. \textit{Phys.\ Rev.\ D}, 34, 470.
\bibitem{Joos1985} Joos, E., and Zeh, H.~D. (1985). The emergence of classical properties. \textit{Z.\ Phys.\ B}, 59, 223.
\bibitem{Bose2017} Bose, S.\ et al.\ (2017). Spin entanglement witness for quantum gravity. \textit{Phys.\ Rev.\ Lett.}, 119, 240401.
\bibitem{Marletto2017} Marletto, C., and Vedral, V. (2017). Gravitationally-induced entanglement. \textit{Phys.\ Rev.\ Lett.}, 119, 240402.
\bibitem{Zurek2003} Zurek, W.~H. (2003). Decoherence, einselection, and the quantum origins of the classical. \textit{Rev.\ Mod.\ Phys.}, 75, 715.
\bibitem{Haag1996} Haag, R. (1996). \textit{Local Quantum Physics}. Springer.
\end{thebibliography}
\end{document}
